\chapter{Meshes, wireframes \& CAD}\label{ch:meshes}
\begin{goals}
\textbf{In this chapter} you open surfaces and solids from the formats mining software writes (Leapfrog, Datamine, DXF, OBJ, PLY, OMF and \file{.wfr}), decide between the exact and the streamed form of a large mesh, style a mesh, and export one to another format.
\end{goals}

A geological model is as much shapes as numbers: the topography, a pit design, a fault, the wireframe of an ore domain. micro opens them as \term{mesh} layers. A mesh layer has its own eye, color and opacity, is cut by sections, and can be picked. The closed ones are solids you can test a block model against (\cref{ch:geometry}). This chapter is about getting meshes in, looking at them, and getting them out.

\section{What opens}\label{sec:meshes-formats}

micro reads the mesh and CAD formats in \cref{tab:meshes-formats}.

\begin{table}[htbp]
\small
\begin{tabularx}{\linewidth}{@{}lL@{}}\toprule
format & what micro reads from it \\ \midrule
\file{.msh} & a Leapfrog (ARANZ) mesh: one surface or solid \\
\file{.lfm} & a Leapfrog model file: several named meshes with their colors \\
\file{.obj} & a Wavefront mesh, with its named objects or groups \\
\file{.ply} & a mesh when the file has faces, with any per-vertex properties as columns \\
Datamine pair & a points file and a triangles file, usually \file{…pt.dm} and \file{…tr.dm}, as one mesh \\
\file{.dxf} & 3D faces and polyface meshes as a mesh; 3D polylines as design strings; points as pegs \\
\file{.wfr} & many surfaces, strings and point sets in one compact file \\
\file{.omf} & the triangulated surfaces of an Open Mining Format project, along with its block models and point sets \\
\bottomrule
\end{tabularx}
\caption{Mesh and CAD formats.}\label{tab:meshes-formats}
\end{table}

Open them like anything else: drag, or \menu{File>Open…} (\cref{ch:opening}). The open dialog describes each file before it loads: \ui{Wavefront mesh (.obj) · 4 objects}, \ui{Leapfrog model (.lfm) · 6 meshes}, \ui{DXF drawing · 12,400 faces → mesh · 86 strings}.

\subsection{Files with several objects}

An \file{.obj} with several named objects, or an \file{.lfm} with several meshes, shows \ui{split into N solids (as a group)} in the open dialog, ticked. Hover it to read the names. Left ticked, each object becomes its own layer, in a group named after the file, and \file{.lfm} meshes keep the colors they had in Leapfrog. Untick it to get one merged mesh.

Split a set of domain wireframes, so that each solid can be shown, hidden and queried by itself. A group of solids is also the input for assigning domains to a model (\cref{ch:geometry}).

\begin{tip}
For pit phases or alternative interpretations, shapes you want to see one at a time, right-click the group and tick \ui{Mutually exclusive}. The group then shows one child at a time, and \key{\{} and \key{\}} step through them. The demo's \ui{mining phases} group is a good place to try it.
\end{tip}

\subsection{Datamine wireframes}\label{sec:meshes-dm}

Datamine stores a triangulated surface as two \file{.dm} files: a points file with the vertices, and a triangles file that refers to them by number. By convention they are named \file{<name>pt.dm} and \file{<name>tr.dm}.

Open both together. The open dialog labels each half, as in \ui{Datamine wireframe · points (pairs with its triangles file)}, and loads each pair as one mesh. micro pairs the files by what their columns say they are, preferring matching base names when you drop several pairs at once. Triangles that refer to a missing point, or that have no area, are left out and counted.

A lone half cannot be drawn. micro says so on the status line and asks for both halves.

\section{Exact and streamed}\label{sec:meshes-tiers}

micro can hold a mesh in two forms. For a large mesh the choice decides whether you can compute with it or only look at it.

\begin{description}
\item[Exact]\index{mesh!exact and streamed} Every vertex and triangle is kept, at full precision. This is what solid queries, mesh QA, rasterizing to a grid and structural analysis need. It costs graphics memory in proportion to the mesh.
\item[Streamed] The triangles are packed into compact chunks and drawn within a budget, like a point cloud. It is light on memory and stays responsive at any size, but it is for looking: micro keeps no geometry to compute with.
\end{description}

A mesh of up to two million triangles opens exact. A larger one opens streamed, so that a 30-million-triangle photogrammetric outcrop does not exhaust the machine before you have seen it. Very large \file{.obj}, \file{.lfm} and \file{.dxf} files likewise open as one streamed mesh, and the open dialog says so in the file's description.

The layer's right-click menu converts between the two with \ui{Convert to streamed} and \ui{Convert to exact}, and states the cost before you commit: each item is followed by the memory the layer uses now and an estimate of what it will use after. Converting a mesh above the cap to exact asks for confirmation, with the triangle count and the memory it will take. There is no hard ceiling. On a workstation, loading a 20-million-triangle mesh exactly is a reasonable thing to want; on a phone it is not.

\begin{note}
Commands that need the exact form are absent from a streamed mesh's menu: \ui{Select by this solid…}, \ui{Rasterize to grid…}, \ui{Extend to footprint…} and \ui{Combine surfaces…}. If you are looking for one of them and cannot find it, the mesh is streamed.
\end{note}

\section{DXF}\label{sec:meshes-dxf}

A DXF file is a drawing, and a drawing can hold several kinds of thing. micro sorts what it finds into three:

\begin{itemize}
\item \textbf{Surfaces.} 3D faces and polyface or polygon meshes become one mesh layer.
\item \textbf{Design strings.} 3D polylines (pit crests and toes, ramps, drive centerlines) become a \term{strings} layer, drawn as thin tubes. Arcs in a polyline are sampled into short chords that lie on the arc. Each string is categorized by its DXF layer, so the \ui{classes} tab of the layer's properties hides the toe lines or one bench with a click (\cref{ch:color}).
\item \textbf{Pegs.} Points become a points layer.
\end{itemize}

A DXF opens as its surfaces when it has any, otherwise as its strings, otherwise as its points. \ui{Open as} in the open dialog offers the other kinds the drawing contains (\ui{mesh (3DFACE surfaces)}, \ui{design strings (capsules)}, \ui{points (pegs)}), and \ui{Reinterpret as…} on the layer crosses between them afterwards.

A drawing too large to preview opens streamed, and its \ui{Open as} offers \ui{everything it holds (streamed)}, which reads all three kinds in one pass, or a single kind. Solids stored as ACIS bodies are not read.

\section{\file{.wfr}: compact wireframes}\label{sec:meshes-wfr}

\file{.wfr}\index{wfr} is an open wireframe format (\url{https://gentropic.org/wireframe}). One file holds many surfaces, strings and point sets, each with its name and attributes, in a fraction of the size of the text formats. micro reads it, writes it, and can use it as the form a project keeps its meshes in.

\subsection{Converting files}

\menu{File>Convert to .wfr…} converts mesh files without loading them as layers. Click \ui{Add files…} and choose \file{.ply}, \file{.obj}, \file{.dxf}, \file{.msh}, \file{.lfm} or \file{.wfr} files, or Datamine pairs, then set:

\begin{table}[htbp]
\small
\begin{tabularx}{\linewidth}{@{}lMD@{}}\toprule
setting & meaning & default \\ \midrule
\ui{output} & the name of the \file{.wfr} to write & — \\
\ui{save to} & \ui{a file on disk (streamed)} (Chrome and Edge), the open project's \file{meshes} folder, or \ui{download} & the project, when one is open \\
\ui{coordinates} & the grid coordinates are stored on: \ui{0.1 mm grid}, \ui{1 mm grid}, \ui{1 cm grid}, or \ui{lossless (float64)} & 0.1 mm grid \\
\ui{memory} & how much the conversion may hold before spilling to disk: 256~MB to 2~GB & 512 MB \\
\ui{verify topology} & check the topology of the result & off \\
\ui{open when done} & load the result as layers & on \\
\bottomrule
\end{tabularx}
\caption{The \ui{convert to .wfr} dialog.}\label{tab:meshes-wfr}
\end{table}

The coordinate grid is the precision of the result. A tenth of a millimeter is finer than any survey; \ui{lossless} keeps the source's numbers bit for bit when you need to prove nothing moved.

\subsection{Storing a project's meshes}

When you save a project that has meshes, strings or pegs opened from other formats, micro offers once to store them in the project as \file{.wfr}: smaller, quick to reopen, attributes kept. The answers are \ui{Never}, \ui{Not now}, \ui{Yes} and \ui{Yes, always}; \menu{Tools>Store meshes as .wfr on save} shows the standing choice and changes it. Source files already in the project are kept: packing never deletes data. The layers keep their identity, names, colors and legend, and their lineage records the store.

\begin{note}
A strings or pegs layer that carries per-record work, such as a painted column or a selection, is left on its source file instead of being packed, because \file{.wfr} renumbers such records and the work would land on the wrong ones. Mesh vertices keep their numbering and are safe.
\end{note}

\section{OMF projects}\label{sec:meshes-omf}

An Open Mining Format file (\file{.omf}, version~1), written by Leapfrog and several other packages, bundles a whole project. The open dialog lists what it holds before you open it, and micro opens one layer per element, in a group named after the file:

\begin{itemize}
\item triangulated surfaces become meshes;
\item point sets become points layers;
\item a block model becomes a block layer with its lattice declared, rotated axes kept, and checked against every block (\cref{sec:blocks-rotated}); a model with uneven block widths opens as centroids;
\item scalar data become numeric columns, and mapped (legend) data a text column.
\end{itemize}

Not read yet, and named on the status line when a file has them: line sets, gridded surfaces, vector and date data, textures, and data attached to surfaces.

\section{Picking, and a mesh's records}\label{sec:meshes-records}

Click a mesh and the record panel reports the point on its surface under the cursor: an elevation read off a topography, a coordinate read off a contact.

A mesh also has records of its own: its \term{vertices}. The attribute table (\key{t}) lists them with their coordinates, and any properties the file stored per vertex arrive as columns. A PLY exported from photogrammetry carries a color per vertex; those are columns of its table, under the names the file gave them (typically \col{red}, \col{green} and \col{blue}), and they can be filtered and exported like any others. OBJ and \file{.msh} files store nothing per vertex, so their tables hold coordinates only.

This holds for a mesh in its exact form. A streamed mesh keeps no records.

\begin{why}
Treating vertices as records looks like bookkeeping, but it is what lets everything micro does to a table (filter, paint, select, export) apply to a surface. Painting a joint plane on an outcrop mesh and measuring its dip (\cref{ch:structure}) is painting a column on the mesh's vertices.
\end{why}

\section{Styling a mesh}\label{sec:meshes-style}

A mesh has one color, its \term[mesh!tint]{tint}, and an opacity.

\ui{Tint…} on the layer's menu offers seven presets chosen to read as geology (\ui{Grey}, \ui{Sand}, \ui{Rust}, \ui{Green}, \ui{Blue}, \ui{Violet}, \ui{Brown}) and \ui{Custom…}, which takes a hex color. The \ui{style} tab of the layer's properties has a color picker for the same thing.

\ui{Opacity} offers 100, 75, 50 and 25\,\% on the menu, and a slider from 5 to 100\,\% in the \ui{style} tab. A translucent mesh is drawn as a fine screen of dots, not blended, so what is behind it is always drawn correctly, whatever the order of the layers. A half-opaque shell drawn whole over a sectioned model reads well.

\textbf{Sections.} A mesh that follows the section is drawn as its \emph{trace}: the line where the plane meets it, laid on the cut face over the blocks, like a wireframe outline on a paper section. The surface itself is not drawn while it follows, so its tint shows only in that line and its opacity shows not at all. \ui{Cut by sections} on the layer's menu sets how this mesh responds: \ui{follow the section}, \ui{follow, with its own slab…}, \ui{keep front only}, \ui{keep behind only}, or \ui{don't cut}. Choose \ui{don't cut} to see the mesh whole, tinted and translucent, over the cut model: a topography for context, or an ore shell around its blocks (\cref{ch:sections}).

\shot{meshes-demo}{Meshes in the demo project over the sectioned deposit. The ore shell is set to \ui{don't cut}, so it is drawn whole, tinted Green at 50\,\% opacity; phase~1 still follows the section and shows only as its trace, the dark line across the cut face.}

A grid shown as a surface is also a mesh layer, but it is colored by elevation or by a draped grid, not by a tint (\cref{sec:grids-surface}).

\begin{tryit}
Style the demo's shapes.
\begin{enumerate}
\item Open the demo project. The ore shell shows only as an orange ring on the cut face: it follows the section, so only its trace is drawn.
\item In the layers panel, right-click \ui{ore shell} and set \ui{Cut by sections} to \ui{don't cut}. The whole shell appears over the cut model.
\item Right-click it again and choose \ui{Tint…}, then \ui{Green}.
\item Right-click it once more and set \ui{Opacity} to \ui{50\%}. The blocks behind it show through.
\item A translucent mesh is picked only while it is the active layer, so click \ui{ore shell} in the layers panel, then click on the shell. The record panel gives the coordinates of the point you clicked.
\item Open the \ui{mining phases} group. Only \ui{phase 1 — starter pit} is shown.
\item Right-click the group and tick \ui{Mutually exclusive}.
\item Click the phase to make it the active layer, then press \key{\}} repeatedly. The group shows each of the three phases in turn.
\item Right-click the phase being shown and set \ui{Cut by sections} to \ui{don't cut}. The line on the cut face becomes the whole pit, drawn whatever the section through the model is doing.
\end{enumerate}
\end{tryit}

\section{What you can do with a mesh}\label{sec:meshes-tools}

The middle of a mesh's right-click menu holds the tools. Each has its own chapter; this is the map.

\begin{table}[htbp]
\small
\begin{tabularx}{\linewidth}{@{}lMD@{}}\toprule
command & what it does & see \\ \midrule
\ui{Select by this solid…} & select, flag or measure the fraction of a model's blocks inside the solid, or above and below an open surface & \cref{ch:geometry} \\
\ui{Rasterize to grid…} & turn the mesh into a gridded surface & \cref{sec:grids-rasterize} \\
\ui{Extend to footprint…} & grow a surface to cover a model's extent & \cref{sec:grids-tools} \\
\ui{Combine surfaces…} & merge two or more surfaces by a rule & \cref{sec:grids-tools} \\
\ui{Mesh QA / volume…} & watertight or not, shells, flipped faces, slivers, volume and tonnage & \cref{ch:qa} \\
\ui{Structural analysis…} & paint planes on an outcrop mesh and read their attitudes & \cref{ch:structure} \\
\ui{Export mesh…} & write the mesh in another format & below \\
\bottomrule
\end{tabularx}
\caption{The tools on a mesh layer's menu.}\label{tab:meshes-tools}
\end{table}

\begin{tip}
Before a wireframe bounds anything, whether a domain flag or a volume in a report, run \ui{Mesh QA / volume…} on it. A solid with a hole gives a volume that looks reasonable and is wrong.
\end{tip}

\section{Exporting meshes}\label{sec:meshes-export}

Right-click a mesh and choose \ui{Export mesh…}. With several mesh layers selected in the panel (\key{⌘}-click), or on a group, the command becomes \ui{Export N meshes…} and writes them together.

\begin{table}[htbp]
\small
\begin{tabularx}{\linewidth}{@{}lL@{}}\toprule
format & several layers become \\ \midrule
\ui{OBJ — Wavefront (text)} & named objects in one file \\
\ui{LFM — Leapfrog model (colours)} & named meshes, each with its layer's tint as its color \\
\ui{PLY — binary, double precision} & one merged mesh \\
\ui{MSH — ARANZ / Leapfrog surface} & one merged mesh \\
\ui{WFR — wireframe (bricked, attributes kept)} & named elements \\
\bottomrule
\end{tabularx}
\caption{Mesh export formats. The dialog notes, for the format you pick, what happens to several layers.}\label{tab:meshes-export}
\end{table}

\shot[0.7\linewidth]{meshes-export-dialog}{\ui{Export N meshes…} on the demo's \ui{mining phases} group, with the LFM format chosen: three layers become three named meshes with their tints.}

The export re-reads the source geometry, so coordinates leave at full double precision. That is why the PLY writer uses doubles: at a UTM northing of seven million meters, single precision would cost about half a meter.

Any mesh micro reads can be written in any of these formats, so micro is also a converter: a Datamine wireframe pair in, an \file{.lfm} out, with the layer names and tints you set.

A grid is exported as a grid, not a mesh: \menu{Surface tools>Export grid…} on a grid layer, as ESRI ASCII, GeoTIFF, Surfer, ESRI binary or XYZ (\cref{sec:grids-export}).
